% Editorial insertion. Complete proofs are in the standalone real-order report.
% Requires the manuscript's standard theorem environments and amsmath.
% Label namespace: realorder:
\section{The sharp real-order threshold and its endpoint atom}
\label{realorder:sec:threshold}

For positive real $a,b$, put $w=a+b$ and
$F_{a,b}(z)=\sum_{n\ge1}H_{n-1}^{(b)}z^n/n^a$.
The following result extends the real-$a\ge1$ signed-kernel theory;
its critical boundary is not represented by an integrable density alone.
Write $k_s(T)=T^{s-1}/\Gamma(s)$ and
\[
 \mathfrak h(t)=\frac1{1-e^{-t}}-\frac1t,\qquad
 C_{a,b}(T)=\frac1{\Gamma(a)\Gamma(b)}
 \int_0^T(T-t)^{a-1}t^{b-1}\mathfrak h(t)\,dt.
\]
Here $\mathfrak h$ increases strictly from $1/2$ to $1$.

\begin{theorem}[Sharp signed-moment threshold]
\label{realorder:thm:threshold}
A finite signed measure $\mu$ on $[0,1]$ satisfying
$\int u^{n-1}\,d\mu=H_{n-1}^{(b)}/n^a$ for all $n\ge1$ exists
if and only if $a+b\ge1$, and it is unique.
When $w>1$, it has density $K_{a,b}(-\log u)$ with
\[
 K_{a,b}(T)=\zeta(b)k_a(T)-\frac{k_{w-1}(T)}{b-1}-C_{a,b}(T)
 \quad(b\ne1),
\]
and
\[
 K_{a,1}(T)=k_a(T)(\gamma+\psi(a)-\log T)-C_{a,1}(T).
\]
This density has exactly one negative--positive crossing in $u$.
When $w=1$, $a=1-b\in(0,1)$ and
\[
 \mu=\frac1a\delta_1-\nu,\qquad
 d\nu(u)=\bigl[-\zeta(b)k_a(-\log u)+C_{a,b}(-\log u)\bigr]du.
\]
The density of $\nu$ is strictly positive and $\nu([0,1])=1/a$.
\end{theorem}

\paragraph{Proof mechanism.}
The regularized Hurwitz integral gives the moments exactly. For $w>1$,
normalizing $K$ by $k_a$ gives, when $b\ne1$,
\[
 \left(\frac{K_{a,b}(T)}{k_a(T)}\right)'
 =-\frac{\Gamma(a)T^{b-2}}{\Gamma(w-1)}
 -\frac{\Gamma(a)T^{b-1}}{\Gamma(w)}
   \mathbb E\bigl[b\mathfrak h(TV)+TV\mathfrak h'(TV)\bigr]<0,
 \quad V\sim\operatorname{Beta}(b,a).
\]
The endpoint signs give a unique crossing; the $b=1$ derivative is
$-1/T-a^{-1}\mathbb E[\mathfrak h(TV)+TV\mathfrak h'(TV)]<0$.
On the critical line the moments tend to $1/a$, forcing an atom of
that mass at one. Below the line the moments are unbounded, precluding
a finite signed measure. Full integrability, continuation and endpoint
arguments are in Sections 2--4 of the accompanying report.

\begin{corollary}[Extended zero geometry]
\label{realorder:cor:zeros}
For $a,b>0$, $a+b\ge1$, the principal branch of $F_{a,b}$ has no
nonzero zeros on $\mathbb C\setminus[1,\infty)$. Its imaginary part
has one simple zero on every upper semicircle $0<\rho\le1$, with
$\arccos\rho<\theta_{a,b}(\rho)<\pi/2$.
\end{corollary}
The usual ordered-sign Stieltjes proof applies. At unit radius the
critical atom requires its own dominance argument; one must not
assume that the negative density is supported away from one.

\begin{theorem}[Critical normalized-radius monotonicity]
\label{realorder:thm:radial}
On $a+b=1$, both orders positive, the function
$\eta(\rho)=\cos\theta_{a,b}(\rho)/\rho$ increases strictly for
$0<\rho\le1$ and belongs to $(1/2,1)$.
Set $x=\rho^2$, $m=2\eta-1$, $d\omega=(1-u)d\nu$ and
$D(u)=1+x(u^2-(1+m)u)$. Then
\[
 \frac{dm}{dx}=
 \frac{\displaystyle\int (u-m)^2(1-u)D(u)^{-2}\,d\omega(u)}
 {(1-xm)\displaystyle\int (1-xu)D(u)^{-2}\,d\omega(u)}>0.
\]
\end{theorem}
This is an additional exact regime, not a proof of the earlier broad
finite-$b$ normalized-radius conjecture.

\paragraph{Euler and boundary scope.}
At $a+b=1$ the ordinary boundary series diverges and boundary values
mean analytic/Abel values. For $a+b\ge1$ the Euler approximants satisfy
$E_N-M2^{-N}<\operatorname{Im}F_{a,b}(i)<E_N$, with $M=1$ for $a\ge1$;
for $0<a<1$ one may take $M=(1-b)^{-1}$ at $b<1$,
$M=b/(b-1)$ at $b>1$, and $M=1+1/a$ at $b=1$.
The report's exact rational certificate proves
$\operatorname{Im}F_{1/10,9/10}(i)<-0.52$, so $M=1$ cannot be
extended without qualification beyond its original $a\ge1$ domain.
